\cotefile{34}
\scanpage{133}
\mspage{V.5}
% archive sheet number: 132
% top diagram
\diagram{\begin{tikzcd}[ampersand replacement=\&,column sep=small,row sep=large]
\& X' \arrow[dl] \& X'_{\unsure{\eta|y}} \arrow[l] \arrow[dl] \& U' \arrow[l] \arrow[dl] \& U'_{\bar\eta} \arrow[l,no head] \arrow[dl] \\
X \arrow[d] \& X_\eta \arrow[l] \arrow[r] \arrow[d,"(\text{pas cart.})"] \& U \arrow[d] \& U_{\bar\eta} \arrow[l,no head] \arrow[d,no head] \\
S \& \eta \arrow[l] \& \eta \arrow[l,equal] \& \bar\eta \arrow[l,no head]
\end{tikzcd}}
% manuscript lines 1-3
1.4. Notons que l'on a
\[
X'=\varprojlim_\alpha X_\alpha,
\]
où les $X_\alpha$ sont étales sur $X$, et affines. On en conclut
\[
U'_{\deleted{\unsure{\bar\eta}}}=\varprojlim_\alpha X_{\alpha\deleted{\unsure{\bar\eta}}}\times_X U,\qquad
U'_{\bar\eta}=\varprojlim_\alpha(X_\alpha\times_X U)_{\bar\eta}
\]
% manuscript lines 4-8
d'où
\[
(\Psi^{(i)}_{\unsure{\nu|v}})_{\bar x}=H^i(U'_{\bar\eta},\Lambda)\simeq
\varinjlim_\alpha H^i((X_{\alpha\deleted{\unsure{\bar\eta}}}\times_X U)_{\bar\eta},\Lambda),
\tag*{(1.4.1)}
\]
où les \deleted{\unsure{$X_{\alpha\bar\eta}$ sont}} $(X_\alpha\times_X U)_{\bar\eta}$ sont des schémas \unsure{algébriques} quasi-affines sur $\bar\eta$, déduits par changement de \unsure{base} des schémas quasi-affines $X_\alpha\times_X U\subset X_\alpha$ sur $\eta$.

% manuscript lines 9-13
Lorsque $U=X_\eta$, cela devient

\deleted{2. Applications}
\[
\deleted{\unsure{\text{alors}}}\quad
(\Psi^{(i)}_\nu)_{\bar x}\unsure{\longrightarrow|=}H^i(\deleted{\unsure{X}}X'_{\bar\eta},\Lambda)\simeq
\varinjlim_\alpha H^i(X_{\alpha\bar\eta},\Lambda),
\tag*{(1.4.2)}
\]
où les $X_{\alpha\bar\eta}$ sont des schémas \inserted{alg.} \uline{affines} sur $\bar\eta$ [déduits des schémas affines $X_{\alpha\eta}$ sur $\eta$]. \deleted{\unsure{les}}

% manuscript lines 14-16
NB. Plus généralement, dans 1.1.1. les $(X_\alpha\times_X U)_{\bar\eta}$ sont affines si on suppose que $U\mathbin{\unsure{-|\to}}X$ est affine, i.e. $U\to X_\eta$ est affine.
